\section{Definition of Discrete Integral}
\label{sec:definition}

The discrete integral accumulates list values into partial sums from a
given initial value. Two representations are equivalent.

\begin{itemize}
  \item Mathematical: $I_k = init + \sum_{i=0}^k L_i$ --- the
    specification
  \item Recursive: $I_0 = L_0 + init$, $I_{k+1} = I_k + L_{k+1}$ --- the
    implementation
\end{itemize}

\subsection{Mathematical Definition}
\label{sec:mathematical-definition}

We define the \textbf{discrete integral} $I = \operatorname{Integral}(L,
init)$ as a list of partial sums such that:

\begin{equation*}
\text{for } k \in [0, n - 1]
\end{equation*}
\begin{equation*}
I_{k} := init + \sum_{i=0}^{k} L_i
\end{equation*}

\subsection{Recursive Definition}
\label{sec:recursive-definition}

The implementation computes the same partial sums by peeling one head from
the list at each recursive step and carrying the current accumulated
value.

\begin{equation*}
\begin{aligned}
I &:= \operatorname{Integral}(L, init) \\
n &:= \lvert L\rvert \\
k &\in [0, n - 1]
\end{aligned}
\end{equation*}

The value of the $k\text{-th}$ element in the integral $I$ is defined
recursively as:

\begin{equation*}
I_k :=
\begin{cases}
L_0 + init & \text{if } k = 0 \\
\operatorname{Integral}(\operatorname{tail}(L),\ \operatorname{head}(L) + init)_{(k - 1)} & \text{if } k > 0
\end{cases}
\end{equation*}

In Scala, this is encoded at
\href{https://github.com/thiagomata/prime-numbers/blob/master/src/main/scala/v1/chapter3/list/integral/Integral.scala}{\texttt{Integral.scala}}:

\begin{lstlisting}[style=scala]
case class Integral(list: List[BigInt], init: BigInt = 0) {
  def apply(position: BigInt): BigInt = {
    require(list.nonEmpty)
    require(position >= 0 && position < list.size)
    if (position == 0) this.head else Integral(list.tail, this.head).apply(position - 1)
  }
  def head: BigInt = {
    require(list.nonEmpty)
    list.head + init
  }
  // ... additional methods omitted
}
\end{lstlisting}
